% Volume IV: Law as an Epistemic Machine
% Part VIII. Legal Friction as Freedom
% Chapter 45: Presumption of Innocence
% Spine: proof-spines/ch045-spine.md

\chapter{Presumption of Innocence}
\label{ch:ch045}

The \textbf{presumption of innocence} is not an empirical prediction that most accused persons are innocent. It is an allocation rule governing who must bear the cost of uncertainty when coercive error is hardest to undo. If \(H_1\) denotes guilt and \(H_0\) innocence, and false conviction and false acquittal carry costs \(C_{FP}\) and \(C_{FN}\), then the decision threshold reflects a normative weighting rather than probability alone:
\[
\text{convict if }
\frac{P(H_1\mid E)}{P(H_0\mid E)}
>
\frac{C_{FP}}{C_{FN}}.
\]
\ToyModel\ The threshold therefore encodes an institutional judgment about asymmetric error costs.
To presume innocence is to forbid accusation from counting as its own confirmation and to assign unresolved uncertainty away from coercive judgment until tested evidence overcomes that default. The presumption changes sign and becomes practical immunity when institutions let that default harden into a refusal to investigate, charge, or revise judgment even after tested evidence has accumulated.
